\section{Revision note}

This version uses geometric torsion, constant complete NY, the physical spherical radius and fixed compact quotients. The epsilon-form action coefficient and spherical normalization of the first edition are corrected, and the vector divergence and endpoint data are treated within one action. Sections 2 and Appendices A--B specify the conventions; the changelog records the changes between editions.

The identification of constant closed-one-form dimensions with physical scalar kernels, the common physical TT spectrum asserted for all four geometries, and the interpretation of the former DBR/RBD difference as a physical ordering preference are withdrawn. Primary preservation requires the multiplier term of (D-01), and quadratic comparisons require matching perturbative order. Geometric closed forms, finite paths and Laplacians are retained in the domains specified in Appendix E.

\stepcounter{equation}\begin{equation}\tag{D-04}\label{eq:D-04}
\begin{aligned}L_{\rm old}&=-8s\dot b+12b^2-12a^2-4as-12\ell a-4\ell s,\\ L_{\rm old,red}&=-\tfrac43\dot s^2+3\ell^2-2\ell s+\tfrac13s^2,\\ \ell&=s/3,\qquad H_{PP}=-3/8.\end{aligned}
\end{equation}

Equation (D-04) is the strictly quadratic trace block of the lowest sphere harmonic in the first edition. Connection elimination and complete primary preservation both give one pair and the same $H_{PP}$. Appendix C gives the scalar comparison derived from the complete action. The supplementary comparison guide supplies the legacy-variable embedding and its relation to the 26-variable block.

This version derives the regular constraint counts for all symmetric homogeneous metrics and the action reduction of nonzero torus modes, connecting the homogeneous cone of Paper06 to compact conditions with TT energy and translation charges; under the hypotheses of Appendix D these conditions are sufficient for local integrability by the Arms--Marsden--Moncrief theorem. It also reduces linear perturbations on selected Nil, Sol and Taub vacuum backgrounds, with quadratic KID conditions and finite-time examples (Appendices F--H). Geometric results used from Papers 03EC, 04 and 06 are attributed and collected in Appendix E. Assumptions and derivations are given in the paper and appendices, with supplement reproduction instructions in the README.
